\section{Related Work}
\label{sec:related}

\citet{dudson2008experiments} studied L-mode discharges using MAST images and BOUT simulations, finding that the mode number of edge turbulence increases with density and decreases with $q_{95}$ and the filament width has the opposite trend. Their simulated filaments are born within about 2~cm of the edge, in the self-generated $E \times B$ shear layer. \citet{benayed2009interelm} looked at filaments in the inter-ELM interval, using both the camera and a reciprocating Langmuir probe, finding that filaments rotate near the LCFS and drift radially outward, that at higher density there are more and faster filaments, and used the radial velocities to test theoretical scaling relations for sheath-connected filaments and inertial filaments.

\citet{kirk2016filaments} scanned the plasma current at fixed density and heating, again in L-mode, and found that the radial velocity and (to a lesser extent) the radial size of filaments decrease with increasing plasma current, consistent with measured divertor profiles. \citet{farley2019identification} did a tomographic inversion of wide-angle images onto equilibrium field lines, recasting the problem of identifying filaments to one of field-aligned identification. All these studies relied on a camera geometry model built with Calcam \citep{silburn2018calcam}.

Our work differs in three ways. First, we don't use wide-angle frames; we use strip mode frames which have only 48 rows of pixels, but are taken at 100~kHz. Second, we derive the pixel scale from the equilibrium model rather than the geometry model. Third, we analyzed hundreds of discharges across four different campaigns; the other studies were limited to a few discharges.

The sheath-connected scaling is due to \citet{krasheninnikov2001scrape}. The inertial scaling is due to \citet{garcia2006fluctuations}. The two-region normalization model used in this paper was introduced by \citet{myra2006collisionality}; the model and its comparison with experiments are reviewed in \citet{dippolito2011convective}. The velocity--size relations were also measured on TORPEX, and the predicted transitions between regimes were observed \citep{theiler2009blob}. 3D simulations of filaments \citep{easy2014three, walkden2015characterization, paruta2018blob} reproduce these scaling relations in form, though realistic geometry and parallel dynamics affect the prefactors. \citet{zweben2007edge} review the imaging of edge turbulence in various devices, including how image intensity relates to density fluctuations.

We used standard detection and tracking methods: connected component labelling after robust thresholding, frame-to-frame assignment according to the Hungarian method \citep{kuhn1955hungarian}, velocity checks using Lucas--Kanade optical flow \citep{lucas1981iterative}, the Theil--Sen estimator \citep{sen1968estimates} for the slope, and the bootstrap \citep{efron1993bootstrap} for confidence intervals. All data came from the FAIR-MAST archive \citep{jackson2024fairmast}.
